\documentclass[a4paper,12pt]{article}

% ------------------------------------------------------------
% Kompatibilny LaTeX zdroj pre pdfLaTeX / Overleaf
% Dokument pouziva iba relativne cesty k obrazkom v priecinku figures/.
% ------------------------------------------------------------
\usepackage[utf8]{inputenc}
\usepackage[T1]{fontenc}
\usepackage[slovak]{babel}
\usepackage{lmodern}
\usepackage{microtype}
\usepackage{geometry}
\geometry{margin=2.5cm}
\usepackage{graphicx}
\usepackage{float}
\usepackage{caption}
\usepackage{booktabs}
\usepackage{array}
\usepackage{longtable}
\usepackage{xcolor}
\usepackage{hyperref}
\usepackage{enumitem}
\usepackage{listings}
\usepackage{tikz}
\usetikzlibrary{arrows.meta, positioning, shapes.geometric}

\hypersetup{colorlinks=true,linkcolor=black,urlcolor=blue,citecolor=black}
\setlength{\parindent}{0pt}
\setlength{\parskip}{6pt}
\setlist[itemize]{topsep=2pt,itemsep=2pt}
\setlist[enumerate]{topsep=2pt,itemsep=2pt}

\lstdefinestyle{terminal}{
    basicstyle=\ttfamily\small,
    breaklines=true,
    frame=single,
    backgroundcolor=\color{gray!8},
    columns=fullflexible
}

\newcommand{\safeimage}[4]{%
\begin{figure}[H]
    \centering
    \IfFileExists{#1}{%
        \includegraphics[width=#2\textwidth]{#1}%
    }{%
        \fbox{\parbox{0.88\textwidth}{\centering\vspace{1.2cm}Obrazok nebol najdeny: \texttt{#1}.\\Pre plnu kompilaciu ponechajte priecinok \texttt{figures/} pri subore \texttt{.tex}.\vspace{1.2cm}}}%
    }
    \caption{#3}
    \label{#4}
\end{figure}
}

\title{\textbf{Technická dokumentácia a používateľská príručka systému zdieľaného supervízneho riadenia autonómnej skupiny dronov vo virtuálnej realite}}
\author{Tímový projekt -- Robotika a kybernetika}
\date{\today}

\begin{document}

\begin{titlepage}
    \centering
    \vspace*{2.2cm}
    {\Large Slovenská technická univerzita v Bratislave\par}
    {\large Fakulta elektrotechniky a informatiky\par}
    \vspace{1.7cm}
    {\Huge\bfseries Príloha 1 - používateľská príručka\par}
    \vspace{0.5cm}
    {\LARGE Zdieľané supervízne riadenie autonómnej skupiny dronov vo virtuálnej realite\par}
    \vspace{1.8cm}
    \vfill
    {\large \today\par}
\end{titlepage}

\tableofcontents
\newpage



\section{Ǔvod}

Tento dokument dopĺňa projektovú dokumentáciu o praktickú používateľskú príručku pre správne nastavenie prostredia a spustenie používateľskej aplikácie. Obsahuje zoznam potrebného hardvéru, inštalačný manuál a praktický návod na spustenie prostredia. Slúži ako stand-alone príručka, ktorú by mal vedieť užívateĺ využiť aj bez znalosti technickej dokumentácie.

\section{Zoznam použitého hardvéru}

Na správne fungovanie aplikácie potrebujete:
\begin{itemize}
    \item Počítač s operačným systémom Ubuntu verzie 22.04 na ktorom bude bežať ROS2 backendová aplikácia a v prípade simulačnej konfugurácie aj simulácia
    \item Počítač s operačným systémom Windows na ktorom bude bežať používateľkská aplikácia a bude cez SteamVR zdieľať obrazovku s VR headsetom
    \item VR headset Pico Neo 3
    \item Dron Airvolute EduDrone 
\end{itemize}


\section{Príprava Windows prostredia}
Windows aplikácia na svoj beh využíva \texttt{MSYS2 UCRT64}, pretože potrebuje kombináciu knižníc PyQt5, GStreamer a Python balíkov, ktoré sa v tomto prostredí inštalujú konzistentnejšie ako cez bežný systémový Python.

Na Windows počítači je potrebné najskôr nainštalovať MSYS2. Po inštalácii sa spúšťa terminál \texttt{MSYS2 UCRT64}; na samotné spúšťanie aplikácie sa nepoužíva klasický Windows \texttt{cmd} ani PowerShell.

\begin{enumerate}
    \item Nainštalujte MSYS2 z oficiálnej stránky projektu: https://www.msys2.org/
    \item Spustite \texttt{MSYS2 UCRT64}.
    \item Aktualizujte systém balíkov:
\end{enumerate}

\begin{lstlisting}[style=terminal]
pacman -Syu
\end{lstlisting}

Ak terminál po aktualizácii vyžiada zatvorenie, zatvorte ho, znovu otvorte \texttt{MSYS2 UCRT64} a príkaz \texttt{pacman -Syu} spustite ešte raz.

\subsection{Inštalácia závislostí aplikácie}

Následne v termínály nainštalujeme všetky
\begin{lstlisting}[style=terminal]
pacman -S --needed \
  mingw-w64-ucrt-x86_64-python \
  mingw-w64-ucrt-x86_64-python-pip \
  mingw-w64-ucrt-x86_64-python-numpy \
  mingw-w64-ucrt-x86_64-python-pyqt5 \
  mingw-w64-ucrt-x86_64-python-gobject \
  mingw-w64-ucrt-x86_64-python-cffi \
  mingw-w64-ucrt-x86_64-python-cryptography \
  mingw-w64-ucrt-x86_64-gcc \
  mingw-w64-ucrt-x86_64-gstreamer \
  mingw-w64-ucrt-x86_64-gst-python \
  mingw-w64-ucrt-x86_64-gst-plugins-base \
  mingw-w64-ucrt-x86_64-gst-plugins-good \
  mingw-w64-ucrt-x86_64-gst-plugins-bad \
  mingw-w64-ucrt-x86_64-gst-plugins-ugly \
  mingw-w64-ucrt-x86_64-gst-libav
\end{lstlisting}

Následne doplňte Python balík \texttt{roslibpy}, ktorý umožňuje komunikáciu aplikácie s ROS bridge websocketom:

\begin{lstlisting}[style=terminal]
python -m pip install roslibpy --break-system-packages
\end{lstlisting}

Pre spoľahlivejšiu prevádzku v prostredí MSYS2 UCRT64 sa používa pevne určená verzia balíka \texttt{autobahn}:

\begin{lstlisting}[style=terminal]
python -m pip install --break-system-packages --force-reinstall --no-deps "autobahn==23.6.2"
\end{lstlisting}

\subsection{Povolenie video portov vo Windows firewalle}

Video streamy z backendu sú prijímané cez UDP porty. Ak ich firewall blokuje, aplikácia sa môže spustiť, ale panely dronov nebudú zobrazovať obraz. PowerShell spustite ako administrátor a povoľte vstupné UDP porty:

\begin{lstlisting}[style=terminal]
New-NetFirewallRule -DisplayName "DroneApp GStreamer UDP" \
  -Direction Inbound -Protocol UDP -LocalPort 2223-2225 -Action Allow
\end{lstlisting}

Použité porty musia byť zosúladené s konfiguráciou GStreamer streamov v ROS~2 backendu. Ak sa v projekte používajú iné porty, je potrebné upraviť aj rozsah vo firewall pravidle.

\section{Spustenie backendu na Ubuntu}

Pred spustením backendu je potrebné zistiť IP adresu Windows počítača, na ktorý sa majú posielať GStreamer streamy. Na Windows počítači použite príkaz:

\begin{lstlisting}[style=terminal]
ipconfig
\end{lstlisting}

Na Ubuntu počítači následne spustite projekt s parametrom \texttt{mode (sim/real)} ktorý určuje, či chceme používať simulačné, alebo reálne drony a  \texttt{--gst-host}, kde zadajte cieľovú IP adresu WINDOWS\_IP pre video streamy:

\begin{lstlisting}[style=terminal]
cd ~/TP-zdielane-supervizne-riadenie
./start_swarm.sh  --mode [sim, real] --gst-host WINDOWS_IP 
\end{lstlisting}

Ďalšie nepovinné parametre sú:
\begin{itemize}
    \item --build - automaticky vykoná colcon build
    \item --drones (1,2,3) - počet dronov, ktoré chceme ovládať
    \item --drone1-fcu - porty jednotlivých dronov (keď sa líšia od preddefinovaných)
\end{itemize}
Príklady:

\begin{lstlisting}[style=terminal]
./start_swarm.sh --build --mode sim --drones 3
./start_swarm.sh --mode sim --drones 1
./start_swarm.sh --build --mode real --drones 1 --drone1-fcu "/dev/ttyACM0:57600"
./start_swarm.sh --mode real --drones 1 --drone1-fcu "udp://:14550@"
\end{lstlisting}

Po spustení backendu by mali bežať ROS2 nody, rosbridge websocket, logika misií a odosielanie GStream video streamov smerom na Windows počítač.

\section{Spustenie aplikácie DroneApp na Windows}

V termináli \texttt{MSYS2 UCRT64} prejdite do adresára aplikácie \texttt{DroneApp}. Cestu je potrebné upraviť podľa reálneho umiestnenia repozitára na počítači:

\begin{lstlisting}[style=terminal]
cd /c/Users/marek/Documents/2.semester/Timovy\ projekt/TP-zdielane-supervizne-riadenie/DroneApp
\end{lstlisting}

Aplikácia sa spúšťa s parametrom \texttt{--ros-host}, ktorý určuje IP adresu Ubuntu počítača, kde beží rosbridge:

\begin{lstlisting}[style=terminal]
python DroneApp.py --ros-host UBUNTU_IP
\end{lstlisting}

Príklad:

\begin{lstlisting}[style=terminal]
python DroneApp.py --ros-host 192.168.0.190
\end{lstlisting}

Po spustení by sa malo zobraziť okno \ref{fig:sim-pic}. V ňom sú panely jednotlivých dronov, telemetrické údaje, stav pripojenia, batéria a video obraz.

\safeimage{img/sim-pic.png}{0.92}{Ukážka zo spustenej aplikácie}{fig:sim-pic}


\subsection{Základná kontrola po spustení}

Po spustení systému odporúčame vykonať kontrola, že Windows počítač vidí rosbridge na Ubuntu:

\begin{lstlisting}[style=terminal]
Test-NetConnection UBUNTU_IP -Port 9090
\end{lstlisting}


\section{Pripojenie na VR rozhranie}

Na pripojenie na VR rozhranie použijete Windows počítač. Potrebujete mať na ňom nainštalované: 

\begin{itemize}
    \item Steam so SteamVR
    \item ALVR launcher - stihneme z oficiálneho Githubu: https://github.com/alvr-org/ALVR/releases/tag/v20.14.1 a nainštalujeme
\end{itemize}

Pre pripojenie je potebné mať obidva zariadenia pripojené na rovnakej sieti. 

Na VR aplikácii si medzi aplikáciami nájdete ALVR client a ten spustíte, v heasete sa vám objavia údaje na zariadenia.  

Následne spustíme Steam a ALVR laucher 
V ALVR launcher cez možnosť Trusted Wireless Devices -> Add device manually pridáte podľa parametrov z headsetu pripojenie. 

Po pridaní a spustení cez Launch Stream VR sa vám v headsete zobrazí stream z obrazovky počítača a aplikácia, ktorú je možné ovládať cez controllery VR zariadenia. 

\section{Ovládanie aplikácie}

Systém využíva VR ovládače s nasledujúcou konfiguráciou vstupov:

\begin{itemize}
    \item \textbf{Pravý ovládač}
    \begin{itemize}
        \item Joystick os Y -- pohyb dopredu/dozadu (\texttt{lin\_x})
        \item Joystick os X -- pohyb vľavo/vpravo (\texttt{lin\_y})
        \item Tlačidlo A -- prevzatie alebo uvoľnenie riadenia aktuálne zvoleného drona
        \item Tlačidlo B -- prepnutie na nasledujúci dron
        \item Grip tlačidlo -- prepnutie na predchádzajúci dron
    \end{itemize}

    \item \textbf{Ľavý ovládač}
    \begin{itemize}
        \item Joystick os Y -- pohyb hore/dole (\texttt{lin\_z})
        \item Joystick os X -- rotácia drona okolo osi Z -- yaw (\texttt{ang\_z})
        \item Grip tlačidlo -- núdzové zastavenie drona
    \end{itemize}
\end{itemize}

Na elimináciu nechcených pohybov joystickov je použitá mŕtva zóna (\textit{deadzone}) s hodnotou 0.20.

\section{Ukončenie aplikácie}
Pri ukončení testu najskôr ukončite manuálny režim, ak je aktívny, a až potom zatvorte desktopovú aplikáciu. Backend na Ubuntu ukončite pomocou skriptu: 

\begin{lstlisting}[style=terminal]
cd ~/TP-zdielane-supervizne-riadenie
./stop_swarm.sh
\end{lstlisting}


\end{document}